\documentclass[tikz,border=3pt]{standalone}
\usepackage[UTF8]{ctex}
\usepackage{amsmath}
\usetikzlibrary{arrows.meta,calc,positioning}

\begin{document}
\begin{tikzpicture}[
  font=\sffamily,
  box/.style={
    draw=black,
    line width=0.55pt,
    align=center,
    minimum width=3.05cm,
    minimum height=0.92cm,
    inner sep=3pt
  },
  titlebox/.style={
    box,
    font=\sffamily\bfseries
  },
  note/.style={
    draw=black,
    line width=0.55pt,
    align=center,
    minimum width=12.6cm,
    minimum height=0.58cm,
    inner sep=3pt
  },
  arrow/.style={
    -{Stealth[length=2.1mm,width=1.4mm]},
    line width=0.55pt,
    draw=black
  }
]

\node[font=\sffamily\bfseries\large] at (5.7,4.45) {分层标定中的装置环节与参数传递};

\node[titlebox] (top1) at (0,3.45) {源端驱动与腔发射};
\node[box] (mid1) at (0,2.15) {单臂波包\\峰位、宽度、尾部};
\node[box] (low1) at (0,0.85) {$\Omega(t)$\\$g$、$\kappa$};

\node[titlebox] (top2) at (3.8,3.45) {时间窗与到达统计};
\node[box] (mid2) at (3.8,2.15) {绝对点击时序\\$\Delta t$ 直方图};
\node[box] (low2) at (3.8,0.85) {$\Delta t_w$\\符合口径};

\node[titlebox] (top3) at (7.6,3.45) {两光子干涉};
\node[box] (mid3) at (7.6,2.15) {HOM 谷深、宽度\\中心偏移};
\node[box] (low3) at (7.6,0.85) {$v_{\mathrm{HOM}}$\\$\sigma_\phi$、$\tau_0$};

\node[titlebox] (top4) at (11.4,3.45) {BSM 与条件态};
\node[box] (mid4) at (11.4,2.15) {模式组成、$F_t$\\CHSH、速率};
\node[box] (low4) at (11.4,0.85) {背景/暗计数\\残余混合、工作点};

\foreach \i in {1,2,3,4} {
  \draw[arrow] (top\i.south) -- (mid\i.north);
  \draw[arrow] (mid\i.south) -- (low\i.north);
}

\draw[arrow] (top1.east) -- (top2.west);
\draw[arrow] (top2.east) -- (top3.west);
\draw[arrow] (top3.east) -- (top4.west);

\node[note] at (5.7,-0.15) {层间传递规则：继承已确定参数与不确定度，固定进入后续拟合。};

\end{tikzpicture}
\end{document}
